\documentclass[11pt,a4paper]{ltjsarticle}

\usepackage{luatexja}
\usepackage{luatexja-fontspec}
\usepackage{lmodern}
\usepackage[top=25mm,bottom=25mm,left=22mm,right=22mm]{geometry}
\usepackage{hyperref}

\hypersetup{
  hidelinks,
  pdftitle={ガウス過程回帰を用いた枚葉露光量制御による線幅ばらつき改善（概要・要旨）}
}

\setlength{\parindent}{1\zw}
\setlength{\parskip}{0.25\baselineskip}

\begin{document}
\thispagestyle{empty}

\begin{center}
  {\Large\bfseries 概要（要旨）}\par
  \vspace{1.5\baselineskip}
  {\large ガウス過程回帰を用いた枚葉露光量制御による線幅ばらつき改善}
\end{center}

\vspace{1.5\baselineskip}

半導体リソグラフィー工程では、レジスト下層の膜厚および屈折率の変動によって露光光の反射・干渉条件が変化し、同一の設定露光量であっても形成線幅が変動する。対象工程の積層構造は、上層側からレジスト、上層SiO$_2$、SiN、下層SiO$_2$およびウェハ基板で構成される。露光光は各膜界面で反射し、各膜の光学膜厚に応じた位相変化を受けた後、レジスト内部で干渉する。

この現象は、微細構造による光の反射・干渉によって特定波長が強められる構造色と同様の光学原理で理解できる。リソグラフィー工程では露光波長が固定されているため、その波長におけるレジスト内電場強度が極大または極小となる近傍を利用することが、膜厚変動に対する基本的なロバスト設計となる。

しかし、SiN膜厚、レジスト膜厚および各膜の屈折率変動に対して、同時に低感度となる積層膜条件を設定することは困難である。対象製品では、SiN膜厚の想定工程変動範囲がSwing curve上の線幅感度の高い領域に位置していた。

対象工程では、従来から最小二乗法（OLS）による線形回帰モデルを用いた枚葉フィードフォワード露光量制御を適用していた。既存モデルはレジスト膜厚を説明変数としておらず、積層膜干渉に由来する非線形な線幅応答を十分に表現できない場合があった。

本検討では、既存の枚葉制御構成を維持しながら、線幅予測モデルをガウス過程回帰（GPR）へ変更した。学習時には実測レジスト膜厚を説明変数に追加し、露光量算出時には製品ごとのレジスト膜厚狙い値を入力した。150枚のモデリングデータを用いて既存OLSと提案GPRのモデル性能を比較検証し、学習に使用していない独立した16枚のウェハで制御効果を評価した。さらに、その同じ16枚を後続工程までそのまま流動し、イメージセンサー特性を確認した。その結果、既存制御と比較して線幅ばらつきが低減し、後続工程後のイメージセンサー特性ばらつきも改善した。これらは実験評価の結果であり、提案GPR制御の量産適用実績はまだない。改善率等の実数値は未提示のため、本稿ではプレースホルダーとした。

\end{document}
